\chapter*{Summary for the Reader}
\addcontentsline{toc}{chapter}{Summary for the Reader}

This page answers the questions the project set out to answer, in one line each, and says where in the report the evidence sits. Every number below appears in a figure or table of the chapter named, and every figure carries a source line naming the files and manifest ids it was built from.

\begin{table}[H]
\centering\footnotesize
\begin{tabular}{p{4.0cm}p{8.4cm}p{2.6cm}}
\toprule
Question & Answer & Where \\
\midrule
Where does the grid stand today? & CAISO demand has been flat at 216 to 224~TWh a year since 2019, but its shape changed: spring midday net load fell from 10.6 to 6.4~GW while the 20:00 net load stayed near 22 to 25~GW; the record peak is 52,061~MW (September 6, 2022); prices at the SCE point were negative in 877 hours of 2025; the accounting carbon intensity fell from 257 to 186~g/kWh. & Chapter~\ref{ch:current}, Figures 2.1 to 2.7 \\
How much of the state's power do data centers use? & About 3 to 4 percent: 7.0 to 9.3~TWh a year against 240~TWh of retail sales, roughly 1,000~MW of peak; a range of estimates, not one number. & Chapter~\ref{ch:current}, Figure 2.8 \\
How high is data center construction? & Nationally from USD 1.6 to 75.2 billion a year (2014 to mid-2026) with a break in early 2022 and 61 percent a year after it; California's own indicators (employment, Silicon Valley construction) did not move, but requests to connect did. & Chapter~\ref{ch:growth}, Figures 3.1 to 3.3 \\
RQ1: how large is the pipeline, by stage, relative to the grid? & 20,677~MW in California (5,086 signed, 6,987 applications, 8,604 inquiries): 40 percent of the record peak, 21 times existing data center load, 5 times the state's forecast peak growth to 2030; as energy, 181~TWh a year if all ran flat, 84 percent of CAISO's demand. & Chapter~\ref{ch:growth}, Figure 3.4, Table 3.9 \\
How has supply grown? & In-state generation 204 to 205~TWh (2010 to 2025) with gas 53 to 36 percent and solar plus wind 3 to 34 percent; capacity 73 to 107~GW; 65 percent of planned generators since 2016 were built; 6.5~GW retire by 2030. & Chapter~\ref{ch:supply} \\
RQ2: the 2030 gap? & Median 50~TWh and 7.2~GW (5th to 95th percentile 19 to 78~TWh, 4.2 to 10.1~GW), of which data centers are 18~TWh and 2.3~GW; the upper bound of every MW running flat is 202 to 242~TWh and 25 to 30~GW. & Chapter~\ref{ch:together}, Figures 5.1 to 5.3 \\
RQ4: what can flexibility absorb? & 3.8~GW at 0.5 percent curtailment (Duke's method on CAISO load; Duke: 5.0), twice the weighted data center load, half the weighted peak gap, a fifth of the upper bound. & Chapter~\ref{ch:together}, Figure 5.4 \\
RQ3: how do the counting regimes compare? & The same requests are 20,677~MW under California's and Texas's rules, 5,086 under PJM's, about 2,400 at ERCOT's realized rates and about 1,000 as measured load; ERCOT scores highest on timeliness, California on coverage. & Chapter~\ref{ch:together}, Figures 5.5 to 5.6 \\
What is the contribution? & An independent replication and audit of the CEC's forecast method: its parameters reproduce its own numbers; the supply side is discounted the same way; the gap comes with its uncertainty and its drivers named. & Chapter~\ref{ch:discussion} \\
\bottomrule
\end{tabular}
\end{table}

\paragraph{Glossary.} \textbf{CAISO}: the California Independent System Operator, the nonprofit that runs the wholesale market and high-voltage grid for about 80 percent of the state's load; the rest is served by other balancing authorities such as LADWP and SMUD. \textbf{Balancing authority}: the entity that keeps supply and demand matched in a defined area; EIA-930 reports hourly load by balancing authority. \textbf{CEC, IEPR, CED}: the California Energy Commission, its biennial Integrated Energy Policy Report, and the California Energy Demand forecast adopted through it. \textbf{Energization request and tiers}: a customer's request for utility service, sorted by the CEC into inquiries, active applications and signed agreements. \textbf{Confidence level, utilization factor, ramp}: the CEC's three discounts, the probability a request is built, the share of requested capacity actually drawn at peak (67 percent), and the years to full load. \textbf{Net load}: demand minus utility-scale solar and wind, the load the rest of the fleet must serve. \textbf{Curtailment}: solar or wind output turned down because it cannot be used or moved. \textbf{DLAP}: a default load aggregation point, the price node at which a utility's load settles. \textbf{ELCC}: effective load-carrying capability, the share of a resource's nameplate that counts toward the peak. \textbf{PRM}: planning reserve margin, the capacity held above the expected peak. \textbf{ERCOT, PJM}: the Texas grid operator and the Mid-Atlantic regional transmission organization, used here for their published queue and forecast rules. \textbf{Firm and non-firm}: PJM's split between requests backed by a service obligation and those that are not. \textbf{Once-through cooling (OTC)}: coastal gas units whose cooling-water permits set their retirement dates. \textbf{Headroom}: the constant load that could be added before demand exceeds the historical peak more than a stated share of the time, given curtailment. \textbf{Monte Carlo, Sobol index}: repeated random draws of the inputs to get a distribution of the result, and the share of that result's variance attributable to each input.
